\section{Physics-Informed Neural Network Method}
\label{sec:pinn-method}

This section describes how we use the PINN~\cite{raissi2019} used in our study. The Lorenz--1960 equations~\cite{lorenz1960},
their coefficients, and the initial state are defined in Section~II. The model
takes time as input and returns the three state variables. Figure~\ref{fig:pinn-pipeline}
shows the full computation used during training.

\begin{figure}[htbp]
    \centering
    \begin{tikzpicture}[
        node distance=3mm,
        stage/.style={draw, rounded corners, align=center,
            text width=0.86\columnwidth, inner sep=4pt, font=\footnotesize},
        flow/.style={-{Stealth}, thick}
    ]
        \node[stage] (input) {Collocation times $t_i$};
        \node[stage, below=of input] (network)
            {Neural network $\mathcal{N}(t_i;\theta)$\\
             Four hidden layers, 60 tanh neurons each};
        \node[stage, below=of network] (trial)
            {Hard trial solution\\$u_T(t_i)=u_0+g(t_i)\mathcal{N}(t_i;\theta)$};
        \node[stage, below=of trial] (residual)
            {Automatic differentiation and Lorenz equations\\
             $r(t_i)=\frac{d u_T}{dt}(t_i)-f(u_T(t_i))$};
        \node[stage, below=of residual] (loss)
            {Physics loss $\mathcal{L}=\frac{1}{3N_c}\sum_i\|r(t_i)\|_2^2$\\
             Full-batch Adam update of $\theta$};
        \draw[flow] (input) -- (network);
        \draw[flow] (network) -- (trial);
        \draw[flow] (trial) -- (residual);
        \draw[flow] (residual) -- (loss);
    \end{tikzpicture}
    \caption{PINN training path used in the reported experiment. Time is mapped
    to a three-state neural output. The output is placed inside a hard trial
    solution. Automatic differentiation supplies its time derivative. The
    derivative and the Lorenz right-hand side form the residual used by the loss.}
    \label{fig:pinn-pipeline}
\end{figure}
\FloatBarrier

\subsection{Network Architecture}

Let
\begin{equation}
    \mathcal{N}(t_i;\theta):\mathbb{R}\rightarrow\mathbb{R}^{3}
    \label{eq:network-map}
\end{equation}
which denotes the neural network. The function takes the input t which is the scaler time and then returns the output that has three entries for $x$, $y$, and $z$. The network is a fully connected  with multilayer
perceptron. It has four hidden layers. Each hidden layer has 60 neurons and uses
 hyperbolic tangent activation. The output layer is linear. For initialization we use
Xavier-uniform weights and zero biases. So,our architecture becomes $1\text{-}60\text{-}60\text{-}60\text{-}60\text{-}3$. It has 11,283
trainable parameters $\theta$ including all weights and biases.

\subsection{Hard Initial-Condition Enforcement}

The network defined above can approximate arbitrary continuous
functions. However, it does not guaranty that the solution starts from
the correct initial state. To satisfy this constraint exactly,we construct a hard trial solution~\cite{lagaris1998, sukumar2022} that contains the network output. The raw network output $\mathcal{N}(t_i;\theta)$ is not used as the final solution.
 Instead, we define the trial solution as
\begin{equation}
    u_T(t_i)=u_0+g(t_i)\mathcal{N}(t_i;\theta),
    \qquad
    g(t_i)=\frac{t_i-t_0}{t_f-t_0},
    \label{eq:trial-solution}
\end{equation}
Here, $u_0$ is the fixed initial state and $[t_0,t_f]$ is the time interval.
For the experiment, we set  $t_0=0$ and $t_f=1$ which gives $g(t_i)=t_i$. At the initial
time, $g(t_0)=0$. Therefore,
\begin{equation}
    u_T(t_0)=u_0.
    \label{eq:hard-ic}
\end{equation}
Since this equality holds for every \(\theta\), the initial condition is satisfied exactly by construction.
 The reported run does not add an initial-condition
penalty to the loss.

\subsection{Automatic-Differentiation Residual}

Once the initial condition is included in the trial solution, the next step is to ensure that the solution satisfies the governing differential equations at the interior time points. To do this, we define a physics residual and evaluate it using automatic differentiation. Then we use the resulting residual to train the PINN. For a trial solution
$u_T=(x_T,y_T,z_T)$, the residual is
\begin{equation}
    r(t_i)=\frac{d u_T(t_i)}{d t}-f\bigl(u_T(t_i)\bigr).
    \label{eq:physics-residual}
\end{equation}
The right-hand side uses the reduced Lorenz system
\begin{equation}
    f(u_T)=
    \left(c_x y_Tz_T,\;c_y x_Tz_T,\;c_z x_Ty_T\right),
    \label{eq:lorenz-rhs}
\end{equation}
where $(c_x,c_y,c_z)=(-0.10,1.60,-0.75)$ for $k=2$ and $l=1$. PyTorch~\cite{paszke2019}
automatic differentiation computes each time derivative in \eqref{eq:physics-residual}.
It uses the computational graph of the trial solution. This avoids finite
differences and allows the residual loss to update $\theta$ through
back-propagation.

\subsection{Collocation Loss and Optimization}

The residual defined in~\eqref{eq:physics-residual} must be driven
toward zero across the entire time domain. Because the residual
cannot be evaluated at every real number in $[0,1]$, we evaluate it
at a finite set of collocation points and minimize their
mean-squared value.We sample $N_c=3000$ time points with a Latin-hypercube design~\cite{mckay1979} on $[0,1]$.
The points are sampled once and then kept fixed . For every point $t_i$, the
three residual components are calculated. The resulting physics-only objective is
\begin{equation}
    \mathcal{L}(\theta)=
    \frac{1}{3N_c}\sum_{i=1}^{N_c}\left\|r(t_i)\right\|_2^2.
    \label{eq:physics-loss}
\end{equation}
In \eqref{eq:physics-loss}, no values from the numerical reference trajectory enter.
The reference is used after training for evaluation.

The loss $\mathcal{L}(\theta)$ is minimized with respect to $\theta$ using gradient-based optimization. At each iteration, backpropagation computes $\nabla_\theta \mathcal{L}$ and updates the network parameters. The choice of optimizer, learning-rate schedule  and all remaining training hyperparameters are reported in Section~\ref{sec:setup}. Algorithm~\ref{alg:hard-ic-pinn} summarizes the complete procedure of our model.





\begin{algorithm}[tbp]
\caption{Training and evaluation of the hard-IC Lorenz--1960 PINN}
\label{alg:hard-ic-pinn}
\begin{algorithmic}[1]
\Require $k=2$, $l=1$, $u_0=(0.5,0.75,1.0)$, $[t_0,t_f]=[0,1]$
\Require $N_c=3000$, $E=20000$, $\eta_0=10^{-3}$, $\eta_f=10^{-4}$, seed $=0$
\State Compute $(c_x,c_y,c_z)$ from the coefficient formulas.
\State Draw and fix $\{t_i\}_{i=1}^{N_c}$ with Latin-hypercube sampling.
\State Initialize the $1$--$60$--$60$--$60$--$60$--$3$ tanh MLP.
\State Use Xavier-uniform weights and zero biases.
\For{$e=1,\ldots,E$}
    \State $u_T(t_i)\gets u_0+\frac{t_i-t_0}{t_f-t_0}\mathcal{N}(t_i;\theta)$
    \State $\frac{d u_T}{d t}(t_i)\gets\operatorname{AD}(u_T(t_i),t_i)$
    \State $r(t_i)\gets\frac{d u_T}{d t}(t_i)-f(u_T(t_i))$
    \State $\mathcal{L}\gets\frac{1}{3N_c}\sum_i\|r(t_i)\|_2^2$
    \State $\eta_e\gets\eta_0+(\eta_f-\eta_0)e/E$
    \State Update $\theta$ with one full-batch Adam~\cite{kingma2015} step using $\nabla_\theta\mathcal{L}$.
\EndFor
\State Evaluate $u_T$ on a separate 1,001-point uniform grid.
\State Compare with the DOP853~\cite{hairer1993} reference and report errors and $\|r(t)\|_2$.
\end{algorithmic}
\end{algorithm}